% !TEX root = ../../hw07.tex
% Source: author's supplied text in chat, 2026-11-06, Q1 (i)--(iii).
% Preserve the MVT, explicit contraction steps, geometric-tail calculation,
% completeness argument, and uniqueness contradiction.
% Unify the repeated derivative-bound declaration as M. Repair the
% continuity wording and missing index cases locally. The author requested
% restoration of the expanded tail before its sigma notation, 2026-11-06.
\begin{proof}[Solution]
  Part (i). Since $f'$ is bounded, $|f'(x)|\le M$
  for some $M\ge0$
  and every $x\in\R$.
  The function $f$ is differentiable over $\R$, so
  it is continuous there.
  Thus, we apply the mean value theorem over the interval $[y,x]$,
  where $x,y\in\R$ and $x>y$, and choose a $c\in(y,x)$ such that
  \[
    f'(c)=\frac{f(x)-f(y)}{x-y}.
  \]
  Choose $L=M+1>0$ to satisfy $L>0$.
  Since $|f'(c)|\le M<L$, we obtain
  \[
    \begin{aligned}
      |f'(c)|=\left|\frac{f(x)-f(y)}{x-y}\right|
        &\implies \left|\frac{f(x)-f(y)}{x-y}\right|\le L\\
        &\implies |f(x)-f(y)|\le L|x-y|.
    \end{aligned}
  \]
  If $x<y$, interchange the two inputs. If $x=y$,
  both sides are zero. Thus the inequality holds for every
  $x,y\in\R$.

  Part (ii).
  Take any $x_0\in\R$ and define
  $x_{n+1}=f(x_n)$ for $n\ge0$.
  Consider the following, where $x_0,\ldots,x_3\in\R$,
  \[
    \begin{aligned}
      x_3&=f(x_2),\\
      x_2&=f(x_1),\\
      x_1&=f(x_0).
    \end{aligned}
  \]
  Assuming a Lipschitz constant satisfying $0<L<1$ exists, we can
  apply the inequality from part (i).
  For this part, $L$ denotes the constant in the
  contraction hypothesis.
  \[
    \begin{aligned}
      |f(x_2)-f(x_1)|\le L|x_2-x_1|
        &\implies |x_3-x_2|\le L|x_2-x_1|,\\
      |f(x_1)-f(x_0)|\le L|x_1-x_0|
        &\implies |x_2-x_1|\le L|x_1-x_0|\\
        &\implies |x_3-x_2|\le L^2|x_1-x_0|.
    \end{aligned}
  \]
  Repeated application of the contraction inequality yields
  \[
    |x_{n+1}-x_n|\le L^n|x_1-x_0|,
    \qquad n\ge0.
  \]
  Suppose $0<L<1$. Then
  \[
    |x_1-x_0|\ge |x_2-x_1|\ge |x_3-x_2|,
  \]
  which implies that each successive step is
  no larger than the previous one.
  \autoref{fig:hw07:contractive-steps} illustrates
  shrinking steps that alternate around a fixed point.
  \begin{marginfigure}
    \centering
    \input{\analysisroot/figures/interval-styles.tex}
\begin{tikzpicture}[interval diagram,x=3.4cm,y=.6cm]
  \draw[->] (0,3.5)--(1,3.5);
  \node[above=5pt] at (.5,3.5) {$1$};
  \draw[->] (1,2.3)--(.5,2.3);
  \node[above=5pt] at (.75,2.3) {$1/2$};
  \draw[->] (.5,1.1)--(.75,1.1);
  \node[above=5pt] at (.625,1.1) {$1/4$};
  \draw[->] (-.07,0)--(1.1,0) node[right=4pt] {$t$};
  \foreach \a/\lab in {0/x_0,.5/x_2,.75/x_3,1/x_1} {
    \fill (\a,0) circle[radius=1.6pt];
    \node[below=6pt] at (\a,0) {$\lab$};
  }
  \draw[gray,fill=white] (2/3,0) circle[radius=1.6pt];
  \node[above=5pt] at (2/3,0) {$x^*$};
\end{tikzpicture}

    \caption[Shrinking steps in a contraction iteration.]{%
      For $f(t)=1-t/2$ and $x_0=0$, the filled
      points are $x_0,x_1,x_2,x_3$. The open point is the fixed point
      $x^*=2/3$. From top to bottom, the arrows show successive steps
      of lengths $1$, $1/2$, and $1/4$, with alternating directions.}
    \label{fig:hw07:contractive-steps}
  \end{marginfigure}
  It follows that the total distance from start to finish is less
  than or equal to the sum of all steps,
  \[
    |x_3-x_0|\le |x_3-x_2|+\cdots+|x_1-x_0|.
  \]
  Generalizing this for arbitrary $m>n$ gives
  the following calculation. For tails shorter than four steps, retain
  only the increments with indices $n\le k<m$ in the expanded list.
  \[
    \begin{aligned}
      |x_m-x_n|
        &\le |x_m-x_{m-1}|+\cdots+|x_{n+3}-x_{n+2}|\\
        &\qquad+|x_{n+2}-x_{n+1}|+|x_{n+1}-x_n|\\
        &=\sum_{k=n}^{m-1}|x_{k+1}-x_k|\\
        &\le\sum_{k=n}^{m-1}L^k|x_1-x_0|\\
        &=|x_1-x_0|\sum_{k=n}^{m-1}L^k\\
        &=|x_1-x_0|\frac{L^n-L^m}{1-L}\\
        &\le |x_1-x_0|\frac{L^n}{1-L}
          \qquad(L^m>0)\\
        &=\frac{|x_1-x_0|}{1-L}L^n.
    \end{aligned}
  \]
  The final bound is independent of $m$.
  Taking the limit as $n\to\infty$, we have
  \[
    \lim_{n\to\infty}\frac{|x_1-x_0|}{1-L}L^n=0,
    \qquad 0<L<1.
  \]
  Hence, for every $\epsilon>0$, there exists $N\in\N$ such that
  \[
    \frac{|x_1-x_0|}{1-L}L^n<\epsilon
    \qquad\forall n\ge N.
  \]
  Consequently, for every $m>n\ge N$,
  \[
    |x_m-x_n|<\epsilon,
  \]
  establishing $(x_n)$ as a Cauchy sequence over $\R$.
  The same bound for $m<n$ follows by interchanging
  the indices, and $m=n$ gives zero.
  By the Cauchy completeness theorem in $\R$, there exists
  $x^*\in\R$ such that $x_n\to x^*$.
  Finally, since $f$ is continuous and $x_n=f(x_{n-1})$, taking
  $n\to\infty$ gives
  \[
    \begin{aligned}
      x^*&=\lim_{n\to\infty}x_n\\
        &=\lim_{n\to\infty}f(x_{n-1})\\
        &=f\left(\lim_{n\to\infty}x_{n-1}\right)\\
        &=f(x^*).
    \end{aligned}
  \]
  Therefore, $x^*$ is a fixed point of $f$.

  For part (iii), assume there are two distinct fixed points
  $x^*=f(x^*)$ and $y^*=f(y^*)$. Applying the contraction inequality
  gives
  \[
    \begin{aligned}
      |f(x^*)-f(y^*)|
        &\le L|x^*-y^*|,\qquad 0<L<1\\
      \implies |x^*-y^*|
        &\le L|x^*-y^*|.
    \end{aligned}
  \]
  Since $x^*\ne y^*$, we have $|x^*-y^*|>0$, so we can divide both
  sides by $|x^*-y^*|$ to give $L\ge1$, which contradicts $L<1$.
  Therefore, two distinct fixed points cannot exist and the fixed
  point $x^*$ is unique.
\end{proof}
